\chapter{Barcelona as a case study}
\label{ch:use_case_barcelona}

Barcelona, the capital of Catalonia in northeastern Spain, occupies approximately 101 km² and was home to roughly 1.7 million residents in 2024, making it one of Europe’s most densely populated cities. Its municipal boundary is delineated by two important geographic features: the Mediterranean Sea to the southeast and the Collserola mountain range to the northwest. These natural limits not only constrain urban expansion but also shape travel patterns, as coastal flatlands yield to the steep slopes of Collserola rising to nearly 450 m (Figure~\ref{fig:Barcelona_use_case}C). Administratively, Barcelona is divided into 10 districts and 73 neighborhoods (Figure~\ref{fig:Barcelona_use_case}A), which display marked socioeconomic contrasts (Figure~\ref{fig:Barcelona_use_case}B and D). Districts such as Sarrià–Sant Gervasi and Les Corts stand out with the highest average incomes, while Nou Barris and parts of Ciutat Vella record much lower income levels. The Eixample district, although more socioeconomically mixed, concentrates the largest number of amenities and points of interest, making it a major hub of urban activity. 


\begin{figure}[hpbt!]
    \centering
    \includegraphics[width=1\textwidth]{0_plots/paper1/main/barcelona_four_panel_comparison.jpg}
    \caption{
        \textbf{Overview of Barcelona.} \textbf{(A)} Administrative division of Barcelona: 10 districts (colored polygons) and 73 neighborhoods (gray boundaries). \textbf{(B)} Average net income per capita by neighborhood in 2022, with darker colors indicating lower income levels. \textbf{(C)} City's cycling infrastructure, including the bike lanes and the BSS stations of \textit{Bicing}, overlaid on topographic contours showing elevation variations. \textbf{(D)} Number of POIs per neighborhood, with darker colors indicating less POIs. POIs were obtained from OpenStreetMap.}
    \label{fig:Barcelona_use_case}
  \end{figure}

Over the past two decades, Barcelona has actively pursued policies to promote sustainable and active transport. Initiatives such as the superblock program, low-emission zones, and the expansion of the cycling network reflect this agenda. Today, the city offers more than 260 km of dedicated bike lanes and nearly 2,000 km of bike-friendly streets (Figure~A1C), complemented by eight metro lines, six tram lines, and over 100 bus routes with more than 8,000 stops. Beyond the city limits, rail networks (FGC and Rodalies) and interurban buses extend connectivity to other municipalities.

How residents make use of this infrastructure is reflected in the 2024 survey \textit{Enquesta de Mobilitat en Dia Feiner} (EMEF). According to this survey~\cite{EMEF2024}, Barcelona residents made 5.75 million weekday trips in 2024, excluding intensive work-related travel, which equals an average of 3.9 trips per person per day. Overall, active mobility is the leading mode, followed by public transport and private vehicles (Table~\ref{tab:modal_split}). However, patterns differ depending on the type of trip. For internal trips within Barcelona, active mobility dominates, whereas for trips connecting Barcelona with surrounding municipalities, public transport is the main mode. Within active mobility, the vast majority of trips are made on foot, while cycling and personal mobility vehicles (PMV) (e.g., e-scooters, segways, or similar devices) together account for less than 4\% of trips.

\begin{table}[hpbt!]
    \centering
    \caption{\textbf{Modal split of trips in Barcelona based on EMEF 2024 survey.}}
    \label{tab:modal_split}
    \begin{tabular}{lccc}
        \hline
        & \textbf{Total (\%)} & \textbf{Internal (\%)} & \textbf{Connections (\%)} \\
        \hline
        Active mobility & 54.5 & 60 & 5 \\
        \quad Walking & 51.3 & -- & -- \\
        \quad Bicycle & 2.4 & -- & -- \\
        \quad PMV (e-scooters, Segways, others) & 0.8 & -- & -- \\
        Public transport & 28.8 & 28 & 53 \\
        Private vehicles & 16.7 & 12 & 42 \\
        \hline
    \end{tabular}
\end{table}

Within this landscape, \textit{Bicing}--the city's public BSS--plays a complementary role. Launched in 2007, it rapidly expanded to more than 400 stations in 2008. A major renewal in 2019 introduced e-bikes and raised the fleet to about 6,700 units, with e-bikes representing 15\% of the total. Moreover, the number of stations was once more increased to over 500 (Figure~A1C). Since then, around 2,000 m-bikes have been converted to e-bikes and additional units have been added, so that by December 2022 e-bikes made up to 47\% of a fleet of nearly 7,000 bicycles. Building on this trend, the system is now undergoing a further expansion, with 30 new stations added between 2024 and 2025, bringing the total to 557. The fleet has also grown to nearly 8,000 bicycles, of which about 5,000 are electric. This steady growth in supply has been accompanied by growing demand: by the end of 2023, the system had approximately 147,000 subscribers and recorded 18 million trips, according to Barcelona's City Council~\cite{AjuntamentBarcelona2024_Bicing}. Empirical studies further reveal distinct usage patterns~\cite{Grau-Escolano2024}: e-bikes are more intensively used in hilly districts near Collserola, whereas m-bikes dominate in the city’s flatter coastal areas.

Taken together, these features make Barcelona a compelling case study. The city combines compactness and high population density with marked topographic contrasts between flat coastal areas and hilly districts, as well as pronounced socioeconomic diversity. Although active mobility is widespread, cycling remains marginal at only 2.4\% of trips, highlighting clear potential for growth. This aligns with the city’s Urban Mobility Plan 2025–2030, which sets the goal of increasing the combined share of cycling and personal mobility vehicles by nearly 50\% by 2030~\cite{PMU2025}. Coupled with the extensive open data resources, these conditions position Barcelona as a highly suitable real-world environment for evaluating the proposed framework.